\subsection{Duality; reflexive modules}

Recall that from now on the domain A is assumed to be Noetherian and integrally closed and that P(A) (or simply P) denotes the set of prime ideals of A of height 1. Every lattice with respect to A is a finitely generated A-module (no. 1, Corollary to Proposition 1).

Let V be a vector space of finite rank over K, V* its dual and V** its bidual; we shall identify V and V** by means of the canonical mapping \(c_{V}\) (Algebra, Chapter II, § 7, no. 5, Theorem 6). Let M be a lattice of V; recall that the dual A-module M* of M is canonically identified with the dual lattice of M, the set of \(x^{*} \in V^{*}\) such that \(\langle x, x^{*} \rangle \in A\) for all \(x \in M\) ; the bidual A-module M** of M is therefore a lattice of V which contains M. Moreover M*** = M*, for the relation \(M \subset M**\) implies \((\mathbf{M}^{**})^{*} \subset \mathbf{M}^{*}\) and on the other hand \(M^{*} \subset (\mathbf{M}^{*})^{**}\) by the above (cf.~Set Theory, Chapter III, § 1, no. 5, Proposition 2).

If p is a prime ideal, Proposition 5 applied with N = A gives the relation \((\mathbf{M}^{*})_{\mathfrak{p}} = (\mathbf{M}_{\mathfrak{p}})^{*}\) , which justifies the notation \(M_{p}^{*}\) for both terms.

THEOREM 1. If \(\mathbf{M}\) is a lattice of \(\mathbf{V}\) , then \(\mathbf{M}^{*} = \bigcap_{\mathfrak{p} \in \mathbb{P}} \mathbf{M}_{\mathfrak{p}}^{*}\) .

Clearly \(\mathbf{M}^*\) is contained in each of the \(\mathbf{M}_{\mathfrak{p}}^{*}\) . Conversely, suppose that \(x^{*} \in \bigcap_{\mathfrak{p} \in \mathbf{P}} \mathbf{M}_{\mathfrak{p}}^{*}\) ; if \(x \in \mathbf{M}\) , then \(\langle x, x^{*} \rangle \in \bigcap_{\mathfrak{p} \in \mathbf{P}} \mathbf{A}_{\mathfrak{p}}\) and, as \(\mathbf{A} = \bigcap_{\mathfrak{p} \in \mathbf{P}} \mathbf{A}_{\mathfrak{p}}\) (§ 1, no. 6, Theorem 4), \(x^{*} \in \mathbf{M}^{*}\) .

COROLLARY. \(\mathbf{M}^{**} = \bigcap_{p \in P}\mathbf{M}_{p}\)

Theorem 1 applied to \(\mathbf{M}^*\) shows that \(\mathbf{M}^{**} = \bigcap_{\mathfrak{p}\in \mathbb{P}}\mathbf{M}_{\mathfrak{p}}^{**}\) . But as \(\mathbf{A}_{\mathfrak{p}}\) is a principal ideal domain (§ 1, no. 6, Theorem 4), \(\mathbf{M}_{\mathfrak{p}}\) is a finitely generated free \(\mathbf{A}_{\mathfrak{p}}\) -module and hence \(\mathbf{M}_{\mathfrak{p}}^{**}\) is canonically identified with \(\mathbf{M}_{\mathfrak{p}}\) (Algebra, Chapter II, § 2, no. 7, Proposition 14), whence the corollary.

For any lattice M with respect to A, the canonical mapping \(c_{M}: M \to M^{**}\) (Algebra, Chapter II, §2, no. 7) identifies an element \(x \in M\) with itself, for x is the unique element y of \(V = V^{**}\) such that \(\langle x, x^{*} \rangle = \langle y, x^{*} \rangle\) for all \(x^{*} \in M^{*}\) , since \(M^{*}\) generates \(V^{*}\) . We shall say that M is reflexive if \(M^{**} = M\) (loc. cit.). As we have above \(M^{*} = (M^{*})^{**}\) , it is seen that the dual of any lattice M is always reflexive.

Remark (1) Let M be a finitely generated A-module; it is immediate that the dual M* of M, identified with a sub-A-module of \(\operatorname{Hom}_{\mathbf{A}}(\mathbf{M}, \mathbf{K})\) is a lattice of the vector K-space \(\operatorname{Hom}_{\mathbf{A}}(\mathbf{M}, \mathbf{K})\) ; in particular, every finitely generated reflexive A-module is isomorphic to a lattice of a suitable vector K-space.

THEOREM 2. If M is a lattice of V, the following conditions are equivalent:

\begin{enumerate}
\def\labelenumi{(\alph{enumi})}
\item
  M is reflexive.
\item
  \(\mathbf{M} = \bigcap_{\mathfrak{p}\in \mathbb{P}}\mathbf{M}_{\mathfrak{p}}.\)
\item
  \(\operatorname{Ass}(\mathbf{V} / \mathbf{M})\subset \mathbf{P}\)
\end{enumerate}

The equivalence of (a) and (b) follows from the Corollary to Theorem 1. If (b) holds, V/M is canonically identified with a sub-A-module of the product \(\prod_{\mathfrak{p} \in \mathbb{P}} (\mathrm{V} / \mathrm{M}_{\mathfrak{p}})\) ; but in fact, it is contained in the direct sum \(\bigoplus_{\mathfrak{p} \in \mathbb{P}} (\mathrm{V} / \mathrm{M}_{\mathfrak{p}})\) : for if \(\mathbf{L} \subset \mathbf{M}\) is a free lattice and \((e_i)_{1 \leq i \leq n}\) a basis of \(\mathbf{L}\) , each of the coordinates \(x_i\) of a point \(x \in \mathbf{V}\) with respect to \((e_i)\) belongs to \(\mathbf{A}_{\mathfrak{p}}\) except for a finite number of values of \(\mathfrak{p} \in \mathbb{P}\) . The relation \(\mathrm{V} / \mathrm{M} \subset \bigoplus_{\mathfrak{p} \in \mathbb{P}} (\mathrm{V} / \mathrm{M}_{\mathfrak{p}})\) then implies:

\[
\operatorname{Ass} (\mathrm{V} / \mathrm{M}) \subset \bigcap_ {\mathfrak {p} \in \mathrm{P}} \operatorname{Ass} (\mathrm{V} / \mathrm{M} _ {\mathfrak {p}}).
\]

As \(V / M_{p}\) is an \(A_{p}\) -module, an element of \(A - p\) cannot annihilate an element \(\neq 0\) of \(V / M_{p}\) , since the elements of \(A - p\) are invertible in \(A_{p}\) ; the elements of \(\operatorname{Ass}(V / M_{p})\) are therefore contained in \(p\) and are \(\neq 0\) , since \(V / M_{p}\) is a torsion \(A_{p}\) -module; as \(p\) is of height 1, necessarily \(\operatorname{Ass}(V / M_{p}) = \{p\}\) if \(V / M_{p} \neq \{0\}\) and \(\operatorname{Ass}(V / M_{p}) = \varnothing\) if \(V / M_{p} = \{0\}\) ; hence \(\operatorname{Ass}(V / M) \subset P\) .

Finally, if condition (c) holds, then

\[
\operatorname{Ass} (\mathbf {M} ^ {* *} / \mathbf {M}) \subset \operatorname{Ass} (\mathbf {V} / \mathbf {M}) \subset \mathbf {P}.
\]

On the other hand, if \(\mathfrak{p} \in \mathbf{P}\) , then it has been seen in the proof of the Corollary to Theorem 1 that \(\mathbf{M}_{\mathfrak{p}}^{**} = \mathbf{M}_{\mathfrak{p}}\) , whence \(\mathfrak{p} \notin \operatorname{Ass}(\mathbf{M}^{**}/\mathbf{M})\) (Chapter IV, § 1, no. 3, Corollary 1 to Proposition 7). We conclude that \(\operatorname{Ass}(\mathbf{M}^{**}/\mathbf{M}) = \varnothing\) , whence \(\mathbf{M}^{**} = \mathbf{M}\) (Chapter IV, § 1, no. 1, Corollary 1 to Proposition 2).

COROLLARY. Let M, N be two lattices of V with respect to A such that N is reflexive. In order that \(M \subset N\) , it is necessary and sufficient that, for all \(p \in P\) , \(M_{p} \subset N_{p}\) .

The condition is obviously necessary and, if it is fulfilled, then

\[
\bigcap_ {p \in P} M _ {p} \subset \bigcap_ {p \in P} N _ {p} = N.
\]

As \(\mathbf{M} \subset \mathbf{M}^{**} = \bigcap_{\mathfrak{p} \in \mathbb{P}} \mathbf{M}_{\mathfrak{p}}\) , certainly \(\mathbf{M} \subset \mathbf{N}\) .

Examples

\begin{enumerate}
\def\labelenumi{(\arabic{enumi})}
\item
  Every free lattice is reflexive.
\item
  Take \(\mathbf{V} = \mathbf{K}\) . For a fractional ideal \(\mathfrak{a}\) of \(\mathbf{K}\) to be a reflexive lattice, it is necessary and sufficient that it be a divisorial ideal by virtue of criterion (b) of Theorem 2 and §1, no. 4, Propositions 5 and 7.
\item
  Let M be a lattice with respect to A; if S is a multiplicative subset of A not containing 0, Proposition 5 of no. 1 shows that \(\mathrm{S}^{-1}(\mathbf{M}^{*}) = (\mathrm{S}^{-1}\mathrm{M})^{*}\) ; if M is reflexive, \(S^{-1}M\) is therefore a reflexive lattice with respect to \(S^{-1}A\) .
\end{enumerate}

PROPOSITION 6. (i) If \(\mathbf{M}_1\) and \(\mathbf{M}_2\) are reflexive lattices of \(\mathbf{V}\) , so is \(\mathbf{M}_1 \cap \mathbf{M}_2\) .

\begin{enumerate}
\def\labelenumi{(\roman{enumi})}
\setcounter{enumi}{1}
\item
  If \(\mathbf{W}\) is a vector subspace of \(\mathbf{V}\) and \(\mathbf{M}\) is a reflexive lattice of \(\mathbf{V}\) , \(\mathbf{M} \cap \mathbf{W}\) is a reflexive lattice of \(\mathbf{W}\) .
\item
  Let V, W be two vector spaces of finite rank over K and M (resp. N) a lattice of V (resp. W). If N is reflexive, the lattice N: M of \(\operatorname{Hom}_{\mathbb{K}}(\mathrm{V}, \mathrm{W})\) (no. 1, Proposition 3) is reflexive.
\item
  \((\mathbf{M}_1 \cap \mathbf{M}_2)_p = (\mathbf{M}_1)_p \cap (\mathbf{M}_2)_p\) for all \(p \in P\) (Chapter II, § 2, no. 4, Theorem 1). If \(\mathbf{M}_1 = \bigcap_{p \in P} (\mathbf{M}_1)_p\) and \(\mathbf{M}_2 = \bigcap_{p \in P} (\mathbf{M}_2)_p\) , then
\end{enumerate}

\[
\mathrm{M} _ {1} \cap \mathrm{M} _ {2} = \bigcap_ {\mathfrak {p} \in \mathrm{P}} (\mathrm{M} _ {1} \cap \mathrm{M} _ {2}) _ {\mathfrak {p}}
\]

whence the conclusion by virtue of Theorem 2.

\begin{enumerate}
\def\labelenumi{(\roman{enumi})}
\setcounter{enumi}{1}
\tightlist
\item
  Similarly \((\mathbf{M} \cap \mathbf{W})_{\mathfrak{p}} = \mathbf{M}_{\mathfrak{p}} \cap \mathbf{W}_{\mathfrak{p}} = \mathbf{M}_{\mathfrak{p}} \cap \mathbf{W}\) , whence
\end{enumerate}

\[
\mathbf {M} \cap \mathbf {W} = \bigcap_ {\mathfrak {p} \in \mathbf {P}} (\mathbf {M} \cap \mathbf {W}) _ {\mathfrak {p}},
\]

which proves (ii).

\begin{enumerate}
\def\labelenumi{(\roman{enumi})}
\setcounter{enumi}{2}
\tightlist
\item
  As \(\mathbf{M}\) is finitely generated, it follows from no. 1, Proposition 5 that \((\mathbf{N}:\mathbf{M})_{\mathfrak{p}} = \mathbf{N}_{\mathfrak{p}}:\mathbf{M}_{\mathfrak{p}}\) ; moreover, the relation \(\mathbf{N} = \bigcap_{\mathfrak{p}\in P}\mathbf{N}_{\mathfrak{p}}\) implies:
\end{enumerate}

\[
\mathbf {N}: \mathbf {M} = \bigcap_ {\mathfrak {p} \in \mathbf {P}} \left(\mathbf {N} _ {\mathfrak {p}}: \mathbf {M} _ {\mathfrak {p}}\right).
\]

For if \(f \in \bigcap_{\mathfrak{p} \in \mathbb{P}} (\mathbf{N}_{\mathfrak{p}}: \mathbf{M}_{\mathfrak{p}})\) and \(x \in \mathbf{M}\) , then \(f(x) \in \bigcap_{\mathfrak{p} \in \mathbb{P}} \mathbf{N}_{\mathfrak{p}} = \mathbf{N}\) , whence \(f \in \mathbf{N}: \mathbf{M}\) ; this shows that \(\mathbf{N}: \mathbf{M}\) is reflexive.

\textbf{Remarks}\par

\begin{enumerate}
\def\labelenumi{(\arabic{enumi})}
\setcounter{enumi}{1}
\item
  If \(\mathbf{M}_1\) and \(\mathbf{M}_2\) are reflexive lattices of \(\mathbf{V}\) , the lattice \(\mathbf{M}_1 + \mathbf{M}_2\) is not necessarily reflexive (cf.~§ 1, Exercise 2).
\item
  If M is a finitely generated A-module and T its torsion submodule, the dual \(M^{*}\) of M is the same as the dual of M/T, since, for every linear form f on M, the image \(f(T)\) is a torsion submodule of A and hence zero. As M/T is isomorphic to a lattice of a vector space over K, it is seen that the dual of every finitely generated A-module is reflexive.
\end{enumerate}

PROPOSITION 7. Let \(0 \to M \to N \to Q \to 0\) be an exact sequence of A-modules. Suppose that \(N\) is finitely generated and torsion-free.

\begin{enumerate}
\def\labelenumi{(\roman{enumi})}
\item
  If \(\mathbf{M}\) is reflexive, then \(\operatorname{Ass}(\mathbf{Q}) \subset \mathbf{P} \cup \{\{0\}\}\) (in other words, every ideal associated with \(\mathbf{Q}\) , is, either (0), or of height 1).
\item
  Conversely, if N is reflexive and \(\operatorname{Ass}(\mathbf{Q}) \subset \mathbf{P} \cup \{\{0\}\}\) , then M is reflexive.
\end{enumerate}

As A is Noetherian, M is also finitely generated; if we write \(V = M_{(K)}\) , \(W = N_{(K)}\) , M (resp. N) is canonically identified with a lattice of V (resp. W) (no. 1, Proposition 1). Consider the two exact sequences:

\[
0 \rightarrow \mathrm{V} / \mathrm{M} \rightarrow \mathrm{W} / \mathrm{M} \rightarrow \mathrm{W} / \mathrm{V} \rightarrow 0
\]

\[
0 \rightarrow Q \rightarrow W / M \rightarrow W / N \rightarrow 0.
\]

\begin{enumerate}
\def\labelenumi{(\roman{enumi})}
\tightlist
\item
  We deduce (Chapter IV, § 1, no. 1, Proposition 3) that:
\end{enumerate}

\[
\operatorname{Ass} (Q) \subset \operatorname{Ass} (W / M) \subset \operatorname{Ass} (V / M) \cup \operatorname{Ass} (W / V).
\]

If M is reflexive, then \(\operatorname{Ass}(V/M) \subset P\) (Theorem 2); on the other hand, clearly \(\operatorname{Ass}(W/V)\) is, either empty, or reduced to \(\{0\}\) ; whence (i).

\begin{enumerate}
\def\labelenumi{(\roman{enumi})}
\setcounter{enumi}{1}
\tightlist
\item
  Similarly:
\end{enumerate}

\[
\operatorname{Ass} (\mathbf {V} / \mathbf {M}) \subset \operatorname{Ass} (\mathbf {W} / \mathbf {M}) \subset \operatorname{Ass} (\mathbf {Q}) \cup \operatorname{Ass} (\mathbf {W} / \mathbf {N}).
\]

The hypotheses therefore imply that \(\operatorname{Ass}(\mathbf{V} / \mathbf{M}) \subset \mathbf{P} \cup \{\{0\}\}\) . But \(\mathbf{V} / \mathbf{M}\) is a torsion A-module and hence \(\{0\} \notin \operatorname{Ass}(\mathbf{V} / \mathbf{M})\) ; Theorem 2 then shows that \(\mathbf{M}\) is reflexive.

PROPOSITION 8. Let R and S be two commutative rings, \(\rho: \mathrm{R} \to \mathrm{S}\) a ring homomorphism and M a finitely generated R-module. Suppose that R is Noetherian and that S is a flat R-module. Then, if M is reflexive, so is the S-module \(\mathrm{M}_{(\mathrm{S})} = \mathrm{M} \otimes_{\mathrm{R}} \mathrm{S}\) .

We know (Chapter I, § 2, no. 10, Proposition 11) that there exists a canonical isomorphism \(\omega_{\mathbf{M}}\colon (\mathbf{M}^{*})_{(\mathbf{S})}\to (\mathbf{M}_{(\mathbf{S})})^{*}\) , such that

\[
\langle x \otimes 1, \omega_ {\mathrm{M}} (x ^ {*} \otimes 1) \rangle = \rho (\langle x, x ^ {*} \rangle)
\]

for \(x \in \mathbf{M}\) , \(x^{*} \in \mathbf{M}^{*}\) . As \(\mathbf{M}\) is a quotient of a finitely generated free R-module \(\mathbf{L}\) , \(\mathbf{M}^{*}\) is isomorphic to a sub-R-module of the dual \(\mathbf{L}^{*}\) and \(\mathbf{L}^{*}\) is free and finitely generated; since \(\mathbf{R}\) is Noetherian, \(\mathbf{M}^{*}\) is therefore also a finitely generated R-module, whence an isomorphism \(\omega_{\mathbf{M}^{*}}: (\mathbf{M}^{**})_{(\mathbf{S})} \to ((\mathbf{M}^{*})_{(\mathbf{S})})^{*}\) such that

\[
\langle x ^ {*} \otimes 1, \omega_ {M ^ {*}} (x ^ {* *} \otimes 1) \rangle = \rho (\langle x ^ {*}, x ^ {* *} \rangle)
\]

for \(x^{*} \in \mathbf{M}^{*}\) and \(x^{**} \in \mathbf{M}^{**}\) . On the other hand, there is an isomorphism \(^t\omega_{\mathbf{M}}: (\mathbf{M}_{(\mathbf{S})})^{**} \to ((\mathbf{M}^{*})_{(\mathbf{S})})^{*}\) , whence by composition a canonical isomorphism:

\[
\phi = \left(^ {t} \omega_ {\mathrm{M}} ^ {- 1}\right) \circ (\omega_ {\mathrm{M} ^ {*}}): (\mathrm{M} ^ {* *}) _ {(\mathrm{S})} \rightarrow (\mathrm{M} _ {(\mathrm{S})}) ^ {* *}
\]

such that, in the above notation:

\[
\langle \omega_ {M} (x ^ {*} \otimes 1), \phi (x ^ {* *} \otimes 1) \rangle = \rho (\langle x ^ {*}, x ^ {* *} \rangle).\tag{1}
\]

We consider now the canonical homomorphism \(c_{\mathbf{M}}: \mathbf{M} \to \mathbf{M}^{**}\) and show that the composite homomorphism:

\[
\psi : \mathrm{M} _ {(\mathbf {S})} \xrightarrow [ c _ {\mathbf {M}} \otimes 1 ]{} (\mathrm{M} ^ {* *}) _ {(\mathbf {S})} \longrightarrow (\mathrm{M} _ {(\mathbf {S})}) ^ {* *}\tag{2}
\]

is just the canonical homomorphism \(c_{\mathbf{M}_{(\mathbf{S})}}\) . This follows immediately from (1) which gives the relations:

\[
\begin{array}{r l} \langle \omega_ {\mathrm{M}} (x ^ {*} \otimes 1), \psi (x \otimes 1) \rangle = & \rho (\langle x ^ {*}, c _ {\mathrm{M}} (x) \rangle) \\ & = \rho (\langle x, x ^ {*} \rangle) = \langle x \otimes 1, \omega_ {\mathrm{M}} (x ^ {*} \otimes 1) \rangle \end{array}
\]

and from the fact that the elements \(\omega_{\mathbf{M}}(x^{*}\otimes 1)\) generate \((\mathbf{M}_{(\mathbf{S})})^{*}\) . This being so, the hypothesis that \(\mathbf{M}\) is reflexive means that \(c_{\mathbf{M}}\) is bijective, hence so is \(c_{\mathbf{M}}\otimes 1\) and therefore \(\psi = c_{\mathbf{M}_{(\mathbf{S})}}\) is bijective, which shows the proposition.

